\section{Trabalhos Relacionados}

\subsection{Predição da Qualidade da Água}
A predição orientada a dados de indicadores de qualidade da água tem sido amplamente explorada com RNAs, modelos baseados em árvores e arquiteturas recorrentes. Revisões e estudos recentes reportam o uso consistente de variáveis físico-químicas e de contexto geoespacial para prever pH, OD, turbidez e indicadores compostos \cite{yan2024,ahmed2019,aladhaileh2021,pan2025}. Abordagens baseadas em LSTM são comuns onde os registros temporais são densos \cite{liu2019,pyo2023,qi2020}, enquanto métodos híbridos e de comitês melhoram a robustez sob comportamento não linear e regimes de amostragem limitados \cite{wu2022,chadalavada2025,kalaivanan2022,palani2008}.

Além da escolha do modelo, uma distinção importante na literatura é o objetivo operacional: máxima acurácia preditiva em infraestrutura centralizada versus robustez sob sensoriamento ruidoso, amostragem esparsa e restrições de campo, uma distinção crítica para águas interiores, onde o comportamento do modelo é afetado pela heterogeneidade local, pela sazonalidade e pela incerteza de medição \cite{yan2024,lowe2022,rahu2023}.

\subsection{Edge AI e Avaliação Energética}
Em implantações de IoT ambiental, o custo de inferência deve ser avaliado em conjunto com a qualidade preditiva. Trabalhos anteriores destacam a relevância do benchmarking em nível de plataforma e do projeto de modelos ciente do hardware \cite{rosero2020,zulkifli2022,lamrini2024,elshafeiy2023}. Estudos recentes de IA neuromórfica e energeticamente eficiente indicam que modelos orientados a eventos podem ser atraentes para cenários restritos, mas as evidências de implantação em hardware de borda heterogêneo ainda estão em desenvolvimento \cite{aral2024,azadi2025,huijbregts2024,aribe2025future}. Essa lacuna motiva o presente trabalho: uma comparação unificada da qualidade de inferência multimodal offline e do comportamento de energia por inferência em múltiplas plataformas práticas.

A maioria dos estudos anteriores reporta métricas algorítmicas ou métricas de hardware isoladamente. Em contraste, a implantação em campo requer uma visão acoplada de latência, vazão e energia incremental sob transições realistas de repouso/carga. Consequentemente, o benchmarking multiplataforma deve ser interpretado como um resultado de sistemas, e não apenas como um resultado de modelo.

\rev{Posicionar este trabalho frente a estudos recentes de monitoramento ambiental neuromórfico esclarece sua novidade. \citet{aral2024} discutem a promessa da Edge AI neuromórfica para monitoramento rural em nível conceitual, sem implantação em campo; \citet{azadi2025} demonstram inferência neuromórfica embarcada para IoT industrial, mas não para regressão de qualidade da água; \citet{huijbregts2024} tratam da eficiência de SNNs em nível de circuito; e \citet{alshareeda2025} implantam monitoramento de água baseado em ESP32 sem modelos neuromórficos. Até onde sabemos, nenhum estudo anterior combina regressão SNN multimodal de alvos físico-químicos, dados de campo de um lago amazônico e energia por inferência medida em hardware das classes microcontrolador, SBC e smartphone em um único pipeline offline, que é a lacuna aqui abordada.}
